% M16, 11.09.2026: Leserführung und Absatzlogik; Anforderungen und Anker erhalten.
% ============================================================
% §4.3 Beobachtbare Ausführung und explizite Informationsführung
% — v04 (Claude)
%
% v04 (Review-Runde 2 + Tiefen-Checks):
%   ② "gilt nur als erfolgt wenn nachweisbar" → instrumentierte-Pfade-
%     Formulierung + Zustandswirkung (Kollege, berechtigt: Logging =
%     Beobachtungsinstrument; nicht-geloggt ≠ nicht-passiert).
%   ③ E43-03 beobachtungsnah + GOLD-FUND: tool-calls.jsonl zeigt
%     toolStep 1 = fs_write OHNE vorherigen Read, evidence-Feld
%     deklariert die Quellen trotzdem (Deklaration-ohne-Zugriff);
%     Historien-Herkunft als Interpretation gekennzeichnet (Eingabe-
%     kontext nicht als Rohdaten geloggt).
%   ④ E43-04 ENTSCHÄRFT nach Metadaten-Check: Arme phasen-/wochen-
%     versetzt (A=phase2_1/Mai, B=phase2B/13.06., C nicht lokalisiert),
%     Modelle teils verschieden → "konstant gehalten" + "unveränderte
%     Prompts" GESTRICHEN (Protokoll=Plan≠Ausführung — korrigiert auch
%     meine v02-Verteidigung dieses Punkts). Freeze-Klärung notiert.
%   (v03 = E-ID-Klammern eingefügt.)
%
% v02 (Review-Runde 1 — kritisch geprüft):
%   ÜBERNOMMEN: ② "entspricht exakt" → "verdeutlicht … nun anhand des
%     beobachteten Systemverhaltens" (unnötig absolut). ③ Herkunfts-
%     Präzisierung — sein bester Punkt, trifft die Primärquelle besser
%     als v01 (Note: "wahrscheinlich … mit der weitergereichten
%     Historie" → Problem ist der fehlende EXPLIZITE Ausweis, nicht
%     Unerkennbarkeit). Schluss-Vorschlag übernommen (Kap. 6 =
%     Realisierung+Verifikation, Kap. 7 = Wirkung unter definierten
%     Bedingungen — kein Vorgriff "verifiziert jede Einhaltung").
%   ZURÜCKGEWIESEN: ① Anker im Fließtext — kollidiert mit der
%     Autor-Entscheidung (b): Prosa ankerfrei, Anker in der
%     Anhang-Zuordnungstabelle (Blaupause §5.1 lässt "Anhangtabelle"
%     als Ankerform zu); Pauschalverweis kommt in 4.1 Absatz 4. Seine
%     Ankerpflicht-Liste = deckungsgleich mit dem Dateikopf hier ✓.
%   BEHALTEN mit Beleg: ④ "konstant gehaltene Laufbedingungen" +
%     "unveränderte Prompts" — durch das Vergleichsprotokoll GEDECKT
%     (wörtlich: "Einzige Variable: phase2ContextStrategy … Alle
%     anderen Parameter sind über alle Vergleichs-Runs identisch",
%     inkl. run-config-Fixierung; im Audit verifiziert).
%
% Grundlage: Kapitel-4-Blaupause §8 (Absatzfolge 1–8, A1+A2) ·
% genese-dossier P1 + Z0–Z3 (AUDITIERT 31.08.) · Beispiel-Pivot-Muster
% (Prosa-Form mit kursiven Absatz-Leitwörtern).
%
% Audit-Belege hinter den Kernaussagen (NICHT im Text — für die
% Anhang-Ankertabelle; Priorität committeter Bestand):
%   A1-Befund: Phase-1-Läufe im Run-Archiv (Eventlog ohne fs_write/
%     request_approval trotz Toolcall-Ankündigung im Text; zusätzl.
%     Bonus-Fall: Approval-Payload behauptete Prüfung aller Artefakte,
%     gelesen wurde eines). Privatnotiz-Wortlaut: phase01-notes Z. 6.
%   A2-Befund: Run 20260526_155811_f151d2 (Artefakt geschrieben OHNE
%     fs_read auf context.md/Transkript; Diagnose: Historie lieferte
%     Inhalte implizit).
%   A/B/C: Vergleichsprotokoll fixierte EINE Variable (Kontextstrategie),
%     alle übrigen Parameter konstant; feste Bewertungsrubrik; Beleg
%     committet: thesis-evidence/D1-direct-vs-topic (Checklisten-
%     Instrument robust ggü. freiem Review). KORREKTE Zuordnung:
%     A=message_passing · B=artifact_state · C=independent_source_reads
%     (Dossier-Fehler behoben; im Text FUNKTIONAL benannt).
%
% Regeln eingehalten: S1/S2-Aussageweite sprachlich markiert ("in
% dokumentierten Läufen", "explorativer Vergleich … motivierte");
% VERBOTEN-Check: keine Verallgemeinerung "Orchestrierung bestimmt
% generell Verhalten" (§8 unzulässige Verkürzung), keine Mess-Verben,
% keine Klassen-/Dateinamen. A1/A2-Wortlaut = Blaupause §8 wörtlich.
% Rückverweis 2.2.3 eingelöst (dort wurde diese Analyse angekündigt).
% MERKER: "Abschnitt 2.2.3" ggf. auf \ref umstellen, sobald die
%   Kapitel-2.2-Labels vorliegen. Budget: ~1,3 Seiten (Ziel 1,2–1,5).
% ============================================================

% v05 (01.09., Kollegen-Runde 4.3 — kritisch geprüft, ALLE 5 Punkte
% berechtigt): ① Strategievergleich auf Plan-Formulierung zurückgebaut
% („geplant/sollten") — mein „variiert wurde allein" hatte den eigenen
% v04-Audit-Befund (phasenversetzte Arme) im Fließtext wieder kassiert;
% Nicht-Isolations-Satz ergänzt. ② C-Lücke explizit: die Leitfrage hing
% an der NICHT ausgeführten Variante → „konnte nicht beantwortet
% werden". ③ Referenzfehler „Diese Annahme" (Bezugssatz war bei der
% Autor-Überarbeitung entfallen) → „Diese Form der Handlungs-
% beschreibung". ④ Freigabe-Wirkung präzisiert (nur operative
% Wirkungen blockiert; Artefaktablage lief direkt — zsmfassung Z. 25).
% ⑤ A2: „ob beziehungsweise welche Quellen er ZUSÄTZLICH selbst liest"
% (Explizitheit, kein Direktzugriffs-Zwang).
% v-KONKRETISIERUNG (01.09., Konkretisierungs-Runde): M2 Machbarkeits-
% Motiv + M3 konkreter Aufbau (Einzelagent→lineare Kette, Nachrichten-
% historie) im Ausgangslage-Absatz; M4 Freigabeschritt-Erklärung vor
% E43-02; Autor-X-Fragen verarbeitet (Log→Anhang per E2, Grafik→E3a
% verworfen); alle Ergänzungen notes-verifiziert (s. Inline-Fundstellen).
\section{Beobachtbare Ausführung und explizite Informationsführung}
\label{sec:beobachtbare-ausfuehrung-informationsfuehrung}

\emph{Ausgangslage.} Zunächst sollte erprobt werden, ob ein möglichst
selbstständig arbeitender Agent ein Meeting-Transkript verarbeiten und
daraus Frühphasenartefakte erzeugen kann. Der erste Aufbau bestand aus
einem einzelnen Agenten, der die bereitgestellten Quellen über
Dateiwerkzeuge las und seine Ergebnisse über Werkzeugaufrufe ablegte.
Ein späterer Stand verband spezialisierte Agenten unter anderem für Kontext,
Anforderungen, Risiken und offene Fragen in einer linearen Kette.
Zwischenergebnisse wurden über die Nachrichtenhistorie weitergegeben.
Die tatsächlichen Handlungen wurden noch nicht durchgängig technisch
erfasst; ein wesentlicher Teil der sichtbaren Ablaufbeschreibung
stammte daher aus den Modellausgaben selbst.

\emph{Erster formativer Befund.} Die frühen Entwicklungsnotizen
beschreiben angekündigte Schreiboperationen ohne tatsächliche
Ausführung. Ein archivierter Lauf enthält bei registriertem
Schreibwerkzeug keinen Schreibaufruf und keine erzeugten
Zielartefakte. Sein Modell-Rohtext ist nicht erhalten; die genaue
Zuordnung zur frühen Notiz bleibt offen
(vgl. Anhang~\ref{anh:evidenzanker}, E43-01).

Dieser frühe Stand sah bereits eine menschliche Freigabe vor.
Wirkungen über die Artefaktablage hinaus, etwa Übergaben an die
operative Arbeitsumgebung, waren deaktiviert. Der Agent sollte
stattdessen zum Abschluss eine Freigabe anfordern und strukturiert
darlegen, was er erzeugt und geprüft hatte. Ein Mensch hätte auf
Grundlage dieser Selbstauskunft entscheiden sollen. In einem
weiteren Fall behauptete der Agent beim Schreiben der Freigabe-Anfrage,
alle Pflichtartefakte erneut gelesen und geprüft zu haben. Das
Protokoll wies jedoch nur einen erneuten Lesezugriff auf ein
erzeugtes Artefakt aus (E43-02).

Damit wurde die in Abschnitt~\ref{sec:tool-use-function-calling}
eingeführte Trennung zwischen sprachlicher Behauptung,
Werkzeuganforderung und technischer Ausführung praktisch relevant.
Als Konsequenz wurde die Instrumentierung ausgebaut: Ereignis- und
Werkzeugprotokolle machten für die erfassten Ausführungspfade
unterscheidbar, welche Handlungen lediglich beschrieben und welche
tatsächlich registriert wurden.

\emph{Zweiter formativer Befund.} Die erweiterte Beobachtbarkeit
machte ein Problem der Informationsführung sichtbar. In einem Lauf
der linearen Agentenkette erzeugte der Anforderungsagent sein
Ergebnis bereits im ersten protokollierten Werkzeugschritt. Vor
diesem Schritt war kein eigener Quellzugriff protokolliert, obwohl
der Aufruf die vorgesehenen Quellen als verwendete Grundlage
auswies (E43-03). Die Weitergabe über die Nachrichtenhistorie war
eine plausible Erklärung, ließ sich dem Protokoll aber nicht direkt
entnehmen. Auf Ebene des Verarbeitungsschritts blieb somit unklar,
wie die verwendeten Informationen zum Agenten gelangt waren.
Geschärfte Promptvorgaben forderten daraufhin einen eigenen,
belegbaren Quellzugriff. Im ersten Folgelauf entstand das
Anforderungsartefakt weiterhin ohne eigenen Leseaufruf. Im zweiten
las der Agent zwar die Kontextdatei, schrieb aber das geforderte
Artefakt nicht; die Validierung meldete dessen Fehlen (E43-05).

Zur weiteren Klärung wurde ein Vergleich von drei Formen der
Kontextführung geplant. Die erste sollte die Nachrichtenhistorie
weiterreichen, die zweite explizit gespeicherte Zwischenartefakte
über einen geteilten Zustand übergeben. In der dritten sollte jeder
Agent das Transkript selbst erneut lesen. Untersucht werden sollte,
ob ein eigener Quellzugriff etwa die Themenabdeckung verbessert.
Ein modellgestützter Prüfagent bewertete die erzeugten Artefakte
anhand einer festen Checkliste (E43-04).

Ausgeführt und bewertet wurden die ersten beiden Varianten; die
dritte wurde nicht umgesetzt. Die Frage nach dem möglichen Vorteil
eines unabhängigen erneuten Quellzugriffs blieb damit unbeantwortet.
Auch die ausgeführten Varianten entstanden nicht unter vollständig
identischen Laufbedingungen. Ihre Unterschiede lassen sich deshalb
nicht allein der Kontextführung zuschreiben. Die protokollierten
Quellzugriffe unterschieden sich jedoch und gaben Anlass, explizite
Informationsflüsse weiterzuverfolgen. Der geplante Vergleich blieb
unvollständig: Beim Versuch, seine Bewertung verlässlich zu machen,
rückte der Reviewprozess selbst in den Mittelpunkt. Diesen Übergang
behandelt der folgende Abschnitt.

\emph{Entwurfsentscheidung.} Aus den Befunden wurden zwei
Festlegungen für die weitere Entwicklung abgeleitet. Erstens werden
Modelltext, Werkzeuganforderung und Systemwirkung getrennt erfasst.
Für die instrumentierten Pfade belegen registrierte Werkzeug- und
Systemereignisse eine technische Handlung; ist ihre Zustandswirkung
maßgeblich, wird zusätzlich das entstandene Artefakt geprüft.
Zweitens werden Informationswege als typisierte Übergaben mit
deklarierten Quellzugriffen modelliert. Für jeden Schritt ist damit
festgelegt und im Protokoll prüfbar, welche Eingaben er erhält und
ob beziehungsweise welche Quellen er zusätzlich selbst liest. Daraus ergeben sich die
ersten beiden Designanforderungen:

\textbf{A1:} Agentische Aktionen müssen außerhalb der sprachlichen
Selbstauskunft des Modells beobachtbar und prüfbar sein.

\textbf{A2:} Informationsquellen und Übergaben zwischen
Verarbeitungsschritten müssen explizit modelliert werden.

\emph{Einordnung und Grenzen.} Die Befunde stammen aus einzelnen
Läufen früher Entwicklungsstände. Sie erlauben weder
Häufigkeitsaussagen noch eine Übertragung auf andere Aufbauten.
Der unvollständige Strategievergleich motiviert eine Designrichtung,
belegt aber keine allgemeine Überlegenheit einer Kontextführung.
Beobachtbarkeit macht fehlerhafte Handlungen feststellbar, verhindert
sie jedoch nicht. Ebenso garantiert eine explizite Übergabe keine
richtige Interpretation ihrer Inhalte. Die technische Realisierung
und Verifikation der Mechanismen zu A1 und A2 beschreibt
Kapitel~\ref{chap:maf-realisierung}. Ihre ausgewählten Wirkungen werden
unter den in Kapitel~\ref{chap:evaluation} definierten Bedingungen
untersucht.

% Historische Inline-Herkunftsnotizen vor M16; aktueller Befund im M16-Bericht.
% [M2/M3 verifiziert 01.09.: phase01-zsmfassung Z. 5 („technischer Durch-
%  stich … der Agent entscheidet selbst, welche Tools er nutzt"; Tools
%  fs_list/fs_exists/fs_read/fs_write + request_approval, Z. 45–47);
%  Kette: phase02-iteration-notes Z. 5/10 („Context -> Requirements ->
%  Risks -> Architecture -> OpenQuestions", linearer MAF-Workflow);
%  Nachrichtenhistorie: phase02-notes Z. 283/370/377 (message_passing,
%  gen_ai.input.messages)]
% [Fundstelle Autonomie-Ideal: thesis-story/fruehphase/phase01-zsmfassung.md
%  Z. 6 („wenn der Host mehr übernimmt … RPA. Das will ich verhindern")]
% [Autor-X-Fragen 01.09. verarbeitet: Log-Auszug → NICHT im Text, sondern
%  als A.2-Druckauszug beim Freeze (E2-Entscheid; Kopierliste); Aufbau-
%  Beschreibung → Ausgangslage-Absatz oben (M3); Grafik → verworfen (E3a).]
% [Fundstelle E43-01: fruehphase/phase01-iteration-notes.md — Iteration 1
%  (Beobachtung wörtlich, Z. 6); Fix dokumentiert: Iteration 2
%  (UseFunctionInvocation → echte TOOL_CALL-Events); s.a. zsmfassung Z. 185]
% [Freigabe-Wirkung präzisiert 01.09. (Kollegen-Punkt 4, belegt):
%  zsmfassung Z. 25 — „Aktionen (Commit/Push, GitHub Issues) sind in
%  Phase 1 nicht aktiv. Stattdessen erzeugt der Agent eine
%  Approval-Request als Datei." — Artefakt-Schreiben lief direkt,
%  operative Wirkungen waren bis zur Freigabe blockiert.]
% [M4 verifiziert 01.09.: phase01-notes Z. 37 (Contract
%  request_approval(action,payload), Payload mit runId), Z. 748–753
%  (Selbstauskunfts-Felder intent/reason/evidence/rationaleComplete),
%  zsmfassung Z. 25 (echte Aktionen deaktiviert, Approval als Datei —
%  HITL-Vorform); Befund: Z. 1000 + 1053–1056 (Run 150331: nur
%  fs_read requirements.md, Payload „klingt so, als seien alle
%  Artefakte geprüft")]
% [Fundstelle E43-02: fruehphase/phase01-iteration-notes.md — Self-Check-
%  Befund Umfeld Z. ~1000; Run 20260519_150331_b9fbd4 dort benannt]
% [Fundstelle E43-03: fruehphase/phase02-iteration-notes.md — Iteration 10
%  (Stress-Test 155811) + Z. 277 ff. (kein fs_read; „wahrscheinlich …
%  weitergereichte Chat-/Tool-Historie"); Lauf-Verifikation 31.08.:
%  logs/agents/Phase2RequirementsAgent/tool-calls.jsonl (toolStep 1 =
%  fs_write; evidence-Feld deklariert Transkript+context.md)]
% [Fundstelle E43-05: fruehphase/phase02-iteration-notes.md — Iter 11
%  („RequirementsPrompt3 fuer belegbaren Quellenzugriff"), Iter 12
%  (Run 589104: „eigener Quellenzugriff bleibt unbelegt"), Iter 13
%  (RequirementsPrompt4)]
% [Fundstelle E43-04: fruehphase/vergleichsprotokoll.md (Plan, Kontroll-
%  variablen, A/B/C-Werte; §9 Hypothesen — u. a. „C zeigt mehr Kat.3-
%  Treffer als A: Eigener Quellenzugriff → bessere Themenabdeckung" =
%  Leitfrage belegt) + qualitaetsrubrik.md Kopf (Zweck „für den
%  A/B/C-Vergleich"); Prüfagent = Jury/Judge (phase02-notes Iter 23–26,
%  Iter 24: automatischer Post-Run-Trigger, judgeModel-Config);
%  B implementiert: phase02-notes Iter 27 (13.06., MAF-nativer Shared
%  State — Framework-Detail bewusst erst Kap. 6);
%  C nie implementiert: Code-Suche 31.08. + Autor-Bestätigung]
% [Methodische Ehrlichkeit 01.09. (Kollegen-Punkt 1+2, berechtigt):
%  „Kette blieb unverändert; variiert wurde allein…" widersprach dem
%  eigenen v04-Audit (Arme phasen-/wochenversetzt, Modelle teils
%  verschieden — Ein-Faktor-Konstanz galt nur für den PLAN, nicht die
%  Ausführung) → Plan-Formulierung („geplant/sollten") + expliziter
%  Nicht-Isolations-Satz + C-Lücke ausbuchstabiert (Leitfrage hing an
%  der NICHT gebauten Variante).]
% [A2-Präzisierung 01.09. (Kollegen-Punkt, berechtigt): A2 fordert
%  EXPLIZITHEIT des Informationswegs, nicht redundanten Direktzugriff
%  jedes Agenten auf die Originalquelle — der spätere Requirements-
%  Konsument liest Claims, nicht das Rohtranskript.]
